\section{Dettagli implementativi}
\label{sec:details}

\subsection{Partizionamento non-i.i.d.\ del dataset}
Il dataset FEMNIST scaricato in formato LEAF è partizionato per scrittore:
gli scrittori sono raggruppati in blocchi contigui e ogni blocco è assegnato
a un worker, cosicché ogni client riceva gli stili di scrittura e le
distribuzioni di etichette di più utenti distinti, non di un singolo
scrittore. Da ciascuna partizione locale è riservato un 10\% per la
validazione (usata per early stopping e come metrica di accuratezza
riportata). Una frazione separata di scrittori del test set LEAF, mai
assegnata a nessun worker, è riservata come \emph{global test set} condiviso:
valutare tutti i worker sullo stesso insieme di scrittori mai visti permette
di misurare la convergenza \emph{funzionale} del protocollo di comunicazione,
non solo la vicinanza dei pesi. Va detto che un global test set di questo
tipo non è detto sia disponibile in uno scenario reale: resta comunque un
agglomerato di dati privati. In questo lavoro lo si
è reso necessario per uno scopo puramente valutativo, misurare la
convergenza funzionale del sistema durante la sperimentazione. Non è applicato alcun pre-processing o data
augmentation: introdurre trasformazioni casuali (rotazioni, rumore) per worker
altererebbe artificialmente l'eterogeneità naturale degli stili di
scrittura, che è proprio l'oggetto di studio di un sistema FL non-i.i.d.
E' chiaro come in uno scenario reale, i veri proprietari dei dati, potrebbero applicare tecniche di data augmentation localmente,
col fine di combattere la reale natura non iid dei dati. Il sistema implementa
anche un \emph{local test set} opzionale (\texttt{local\_test\_set: true}):
un'ulteriore frazione di campioni riservata all'interno della partizione di
ciascun worker, dagli \emph{stessi} scrittori già presenti in train/val di
quel worker ma con campioni distinti, tenuti fuori dagli aggiornamenti del
gradiente — misura la generalizzazione su campioni mai visti di scrittori
già noti al worker, non su scrittori nuovi. Non è stato usato nella campagna sperimentale di questo
lavoro (Sezione~\ref{sec:results}): il global test set risponde già alla
domanda più rilevante per un sistema FL P2P — la convergenza
\emph{funzionale tra worker} su scrittori ignoti a tutti — mentre il local
test set misurerebbe una generalizzazione più ristretta, sugli stili di
scrittura già visti da un singolo worker.

\subsection{Ciclo di training a tre fasi}
Ogni round di ogni worker è composto da tre fasi, cronometrate
indipendentemente e riportate nella Sezione~\ref{sec:results}.

La sincronizzazione avviene su un numero fisso di passi $H$, non su epoche:
le partizioni non sono bilanciate tra worker (Sezione~\ref{sec:details}A),
quindi un criterio "un round = un'epoca" farebbe comunicare più spesso i
worker con meno dati, rendendo la frequenza di comunicazione una
conseguenza indiretta del partizionamento anziché un parametro controllato
direttamente.

La struttura è a tre livelli annidati, dal più piccolo al più grande. Il
\emph{minibatch} è l'unità elementare: \texttt{batch\_size}$=32$ campioni,
un passo di SGD (\texttt{train\_step}). L'\emph{epoca} è un giro completo
sulla partizione locale del worker — l'ultimo minibatch incompleto viene
scartato (\texttt{drop\_last=True}), per non avere batch di dimensione
variabile che destabilizzerebbero le statistiche di Batch Normalization. Il
\emph{round} — l'unità su cui il sistema sincronizza — è $H$ minibatch
consecutivi: gli \emph{inner step} di DiLoCo (Sezione~\ref{sec:background}B)
sono letteralmente $H$ ripetizioni di un minibatch, non $H$ epoche, e
possono quindi coprire una piccola frazione di un'epoca, un'epoca esatta, o
più epoche a seconda di come $H\times32$ si confronta con la dimensione
della partizione — il caso concreto è discusso in Fase B qui sotto.

\begin{itemize}
  \item \textbf{Fase A — aggregazione e valutazione.} Il worker applica
        FedAvg (Eq.~\ref{eq:fedavg}) tra i propri pesi e la somma pesata
        accumulata nel buffer di ricezione, quindi valuta il modello risultante
        sul proprio validation set (controllando l'early stopping) e, se abilitato, sul global test set. La
        Fase A include quindi non solo la media dei pesi (costo trascurabile)
        ma anche l'inferenza completa su questi due insiemi di dati.
  \item \textbf{Fase B — addestramento locale.} Il worker esegue $H$ passi di
        SGD (AdamW, batch size 32) su un iteratore che ricicla il proprio
        training set indefinitamente, cosicché ogni round processi esattamente
        $H \times \text{batch\_size}$ campioni indipendentemente dalla
        dimensione della partizione locale. L'iteratore è un generatore
        continuo (\texttt{while True: yield from loader}), creato una sola
        volta prima del ciclo dei round: se
        un'epoca termina durante
        i passi $H$ di un round, il generatore prosegue trasparentemente
        nell'epoca successiva con un nuovo shuffle, senza che il confine di
        round interferisca. Il confine tra un'epoca e la successiva può quindi effettivamente cadere
        durante l'esecuzione di un round anziché sul suo limite. Se un
        tale round consuma le ultime $a$ osservazioni dell'epoca in esaurimento
        (di dimensione $E$, la partizione locale del worker) e prosegue leggendo le restanti
        $H\times\text{batch\_size} - a$ osservazioni da un'epoca appena rimescolata,
        una lettura ripetuta \emph{all'interno dello stesso
        round} si verifica solo se uno di questi nuovi campioni coincide con uno dei
        $a$ già letti in questo stesso round. Poiché il nuovo shuffle è uniforme
        sull'intera partizione dei campioni, la probabilità che ciò accada per un
        singolo campione è bassa. È dunque strutturalmente
        più probabile che un campione della nuova epoca corrisponda a una delle
        $E-a$ osservazioni lette nei round precedenti di quella stessa epoca.
  \item \textbf{Fase C — push one-hop.} Il worker seleziona $\min(k,
        |\text{peer disponibili}|)$ vicini \emph{casualmente e in modo
        indipendente a ogni round} dalla propria cache di peer, e invia loro uno snapshot
        dei pesi correnti via gRPC con timeout configurabile. La selezione
        casuale non è un vincolo del protocollo ma un modello neutro di
        qualunque politica di selezione dei vicini si volesse adottare in
        pratica — prossimità di rete, similarità tra i dati locali dei
        worker, o un qualsiasi altro stato osservabile del nodo; ha inoltre il
        vantaggio di non richiedere alcuna conoscenza della topologia
        sottostante. Se un peer non risponde entro
        il timeout, il worker aggiorna la propria cache interrogando di nuovo
        il discovery server — al più una volta per round, solo in caso di
        fallimento — e tenta un singolo invio verso un peer sostitutivo non
        ancora contattato in quel round. Se anche questo secondo tentativo
        fallisce, il worker non ritenta ulteriormente né richiama di nuovo il
        registry in quel round: prosegue senza bloccarsi sul nodo caduto. La
        scelta mantiene il traffico verso
        il discovery server limitato a una sola chiamata per round anche nel
        caso peggiore, accettando che un peer possa restare non contattato
        per quel round.
\end{itemize}

Un controllo di \emph{staleness} lato ricevente scarta i messaggi il cui
round di provenienza è più vecchio di \texttt{max\_staleness} round rispetto
al round corrente del ricevente, evitando che aggiornamenti molto obsoleti
degradino la convergenza. L'early stopping (pazienza configurabile) è una
decisione puramente locale basata sulla validation loss del singolo worker:
un worker può terminare il proprio training mentre altri continuano,
riducendo nel tempo il numero di vicini disponibili per chi prosegue. Lo
stesso criterio determina anche quale checkpoint (\texttt{model\_best.pt}) viene salvato.
Usare la performance sul global test set per scegliere quale round tenere
costituirebbe una forma di \emph{data leakage}, gonfiando artificialmente il risultato riportato
verso il round che per puro rumore statistico va meglio su quel set
specifico. Il global test set resta così una valutazione mai usata
per decisioni, solo per misurare a posteriori.

\subsection{Modello}
Il modello (Tabella~\ref{tab:model}) è una CNN
con due blocchi convoluzionali (32 e 64 canali, batch normalization, dropout
spaziale) e un classificatore fully-connected, per un totale di 1\,704\,350
parametri float32 — circa 6,82\,MB per snapshot serializzato, la quantità di
dati effettivamente trasmessa a ogni push.

\begin{table}[t]
  \centering
  \caption{Conteggio dei parametri della CNN.}
  \label{tab:model}
  \footnotesize
  \begin{tabular}{lr}
    \toprule
    \textbf{Componente} & \textbf{Parametri} \\
    \midrule
    Blocco conv.\ 1 (1$\to$32, 32$\to$32 canali) & 9\,696 \\
    Blocco conv.\ 2 (32$\to$64, 64$\to$64 canali) & 55\,680 \\
    Classificatore (3136$\to$512$\to$62, con BN) & 1\,638\,974 \\
    \midrule
    \textbf{Totale} & \textbf{1\,704\,350} \\
    \bottomrule
  \end{tabular}
\end{table}

La progressione di canali $1\to32\to64$ approssima una conservazione del volume
di informazione a ogni dimezzamento della risoluzione spaziale
($28\to14\to7$); il kernel $3\times3$ in cascata ottiene il
campo ricettivo di una $5\times5$ con meno parametri, e il same padding
preserva la dimensione spaziale fino al pooling. Un'architettura più pesante
(es.\ ResNet-18, $\sim$11M parametri, $\sim$6$\times$ questo modello) abbiamo ipotizzato fosse
sovradimensionata a questa risoluzione. In un sistema FL dove il modello
viene interamente ritrasmesso a ogni push, la dimensione è fondamentale:
non solo come capacità, ma come costo di comunicazione diretto. Il modello qui usato è quindi un compromesso deliberato tra
capacità sufficiente per il compito e leggerezza per il contesto FL.

Per lo stesso motivo è stato escluso anche un Transformer (es.\ Vision
Transformer): a differenza della convoluzione, il self-attention non ha
alcun \emph{inductive bias} di località o traslazione-invarianza, e deve
apprenderlo dai dati — motivo per cui i ViT richiedono pre-training su
dataset enormi per essere competitivi con le CNN. Le partizioni non-i.i.d.\
di poche centinaia di campioni per classe qui disponibili sono l'opposto
di quel regime. A parità di capacità, un Transformer avrebbe anche più
parametri da ritrasmettere a ogni push rispetto a questa CNN, aggravando
proprio il costo di comunicazione che il modello è pensato per minimizzare.

Batch normalization stabilizza il training; nel contesto federato $\gamma,\beta$ sono aggregati da FedAvg.
Dropout spaziale ($p=0{,}25$) e dropout standard ($p=0{,}5$ sul layer più esposto a
overfitting con $\sim$1,6M parametri) regolarizzano le partizioni locali.
L'inizializzazione dei pesi usa il default di PyTorch, Kaiming uniform per
\texttt{Conv2d} e \texttt{Linear} (He initialization). Il modello è addestrato con
AdamW.

\subsection{Tolleranza ai guasti}
Il sistema include meccanismi opzionali di fault injection configurabili
(probabilità di crash simulato via \texttt{sys.exit} e di drop silenzioso di
un messaggio), implementati ma non esercitati in nessuno dei 16 run
riportati: la loro
validazione empirica resta lavoro futuro. Il comportamento previsto dal codice è il seguente: un crash
simulato interrompe il round corrente con \texttt{sys.exit}, che attraversa
comunque il blocco di terminazione che deregistra il worker — un
comportamento più vicino a un arresto pulito che a un crash reale; gli altri
worker non vengono bloccati e si accorgono dell'assenza del peer al primo
fallimento di invio verso di esso, aggiornando la propria cache. Solo un
segnale SIGKILL, non intercettabile, produce un'uscita ungraceful,
che lascia una entry non valida nel discovery server.

\subsection{Motivazione degli iperparametri di training}
Il gradient clipping ($\|\nabla\|\le1$) limita la deriva dei pesi tra
aggregazioni; il label smoothing ($\epsilon=0{,}1$)
attenua l'eccesso di confidenza sulle classi ambigue di FEMNIST (es.\
\texttt{0}/\texttt{O}, \texttt{1}/\texttt{l}/\texttt{I}).

L'ottimizzatore è AdamW anziché SGD o Adam
semplice, per due motivi. Primo, AdamW disaccoppia il weight decay
dall'aggiornamento basato sul gradiente adattivo: in Adam classico, la
regolarizzazione L2 interagisce con i learning rate adattivi per-parametro
in modo che ne attenua l'effetto in modo non uniforme tra i parametri;
AdamW applica il weight decay direttamente ai pesi, indipendentemente dalla
scala del gradiente, un comportamento più prevedibile per un iperparametro
 che qui resta comunque fisso. Secondo, un
ottimizzatore adattivo per-parametro richiede meno tuning del learning rate
di SGD per convergere in modo ragionevole.

Seguono quindi due categorie. \emph{Iperparametri ML fissi}:
learning rate, batch size, dropout, label smoothing, gradient clipping. \emph{Iperparametri di
sistema variati} (Sezione~\ref{sec:results}): $H\in\{100,500,1000\}$, il fanout $k$ e $N\in\{3,5,8\}$.

\subsection{Piattaforma software}
Il sistema è implementato in Python 3.11 con PyTorch (modello e training),
gRPC/Protobuf (comunicazione push tra worker), Flask (discovery server), containerizzato
con Docker e orchestrato con Docker Compose in locale. Il deployment su AWS
usa Terraform per provisionare un'istanza EC2 per worker più una per il
discovery server, con build dell'immagine Docker in parallelo su ciascuna istanza.

Le immagini Docker sono costruite sfruttando il \emph{layer caching}: ogni
istruzione del Dockerfile produce un layer immutabile riutilizzato se il suo
hash coincide con la build precedente, e l'invalidazione è a cascata (un
layer modificato invalida tutti quelli successivi). Il Dockerfile del worker
copia e compila \texttt{proto/gossip.proto} \emph{prima} di copiare il resto
del sorgente Python: il layer più costoso, l'installazione di PyTorch
($\sim$750\,MB), viene rieseguito solo se cambia \texttt{requirements.worker.txt},
non ad ogni modifica al codice — una modifica al codice invalida solo l'ultimo
layer, riducendo la rebuild da minuti a pochi secondi. Tutti i worker di un
deployment condividono inoltre un'unica immagine.

\subsection{Vincolo di RAM e caricamento del dataset}
\texttt{FEMNISTDataset.\_\_init\_\_}
converte l'intero split di training in tensori tenuti in memoria all'avvio,
una scelta di semplicità che però si scontra con il vincolo di RAM delle
istanze \texttt{t3.small} usate su AWS (2\,GB, Sezione~\ref{sec:results}). Il
formato LEAF serializza i pixel come numeri in JSON: durante il parsing, ogni
valore diventa un oggetto \texttt{float} Python, che occupa più byte rispetto un \texttt{float32}
compatto.

La correzione separa il parsing dal caricamento a runtime.
\texttt{split\_dataset.py} esegue ora una passata che rilegge il JSON
già scritto per ciascun worker e lo converte, una sola volta, in array NumPy
tipizzati (\texttt{float32} + \texttt{int64}), salvati come
\texttt{data.npy}/\texttt{labels.npy} accanto al JSON originale.
\texttt{dataset.py} carica questi binari con \texttt{np.load} quando
presenti, con fallback al JSON per gli split generati prima della modifica:
niente più oggetti \texttt{float} Python intermedi, il picco di RAM scende
e \texttt{torch.from\_numpy} condivide la memoria con l'array NumPy
senza copie aggiuntive. Il \emph{global test set}, letto in sola lettura da
tutti i worker, usa inoltre \texttt{np.load(mmap\_mode=\textquotesingle
c\textquotesingle)}: nelle modalità in cui più worker condividono la stessa
macchina (locale, single-EC2) il sistema operativo mappa le stesse pagine
fisiche per tutti i processi che aprono il file, azzerando la duplicazione
che si avrebbe altrimenti per worker; in modalità multi-istanza, dove ogni
EC2 ospita un solo worker, questo beneficio di condivisione non si applica.

\subsection{Riusabilità}
Il sistema è progettato per non dipendere dalle specificità di FEMNIST. Il
partizionamento non-i.i.d.\ e l'aggregazione pesata per numero di campioni
(Eq.~\ref{eq:fedavg}) non assumono client bilanciati o distribuzioni simili
tra loro: l'eterogeneità è il caso di base gestito dal progetto, non
un'eccezione da accomodare — un dataset LEAF diverso (es.\ Shakespeare,
Sentiment140) o una partizione non-i.i.d.\ più estrema richiederebbero solo
un nuovo caricatore dati (\texttt{core/dataset.py}), non una riprogettazione
della logica di training o di comunicazione. Il protocollo di comunicazione
stesso è agnostico rispetto all'architettura del modello — serializza un
intero \texttt{state\_dict} PyTorch senza logica specifica. Infine, il
numero di client, il fanout $k$ e tutti gli iperparametri sono parametrizzati
in un unico file di
configurazione: adattare il sistema a una nuova scala non richiede modifiche
al codice. La scelta di gRPC su Protobuf è motivata anche da questo obiettivo di
riusabilità: lo schema \texttt{gossip.proto} (Sezione~\ref{sec:design}C) è
indipendente dal linguaggio e gRPC genera client/server per numerosi
linguaggi, per cui un worker non Python potrebbe in linea di principio
interoperare in un contesto eterogeneo.
